\section{The projective-pair bound for plain accumulators}\label{app:fresh-pairs}

Plain accumulators have no random coefficients, but through the random
permutations the bound below still charges for some dependent
cancellations. We use the Newton precode $B_8$ and the
prime $p=2^{128}-159$ of Appendix~\ref{app:all-one}. For the same
encoder, the following refinement proves minimum distance at least
$209716$ ($>10\%$), except with probability less than $2^{-47.91}$ over
the two independent uniform permutations. The bound holds simultaneously
for all nonzero messages, including a message chosen after the
permutations are revealed. Appendix~\ref{app:complete-runs} later
refines it to $11.7\%$.

As in Section~\ref{sec:distance-proof}, fix a source support $S$ with
shortened source space $V_S$, and let $W\subseteq V_S^*$ be the span
of zero equations imposed so far. The \emph{active labels} of a scan
are the coordinates of its incoming vector that are required to be
nonzero; the permutation places them in uniformly random order on the
nonzero input positions. Each active label carries a
linear form on $V_S$, its \emph{required active form}, which must
remain nonzero in $V_S^*/W$. At the start of a run, expose its first two
active input labels jointly. Before either label is revealed, they form
a uniform ordered pair from the $m\ge2$ remaining active labels.
If any required active form vanishes in $V_S^*/W$, the branch has no
surviving message and can be discarded. Otherwise every remaining active
form is nonzero. For the moment, draw two labels independently with
replacement. Let $P_\parallel$ be the probability that their forms are proportional,
and let $P_-$ be the probability that they are exact opposites. For each
one-dimensional subspace $\ell$, let $\mu_\ell$ be its total probability
mass, and write $\mu_v$ for the mass of each form $v$. Partition
the nonzero forms of $\ell$ into pairs $\{v,-v\}$. Odd characteristic
makes these pairs disjoint, and their masses satisfy
\[
 2\sum_{\{v,-v\}\subset\ell}\mu_v\mu_{-v}\le \mu_\ell^2/2.
\]
Indeed, orient each opposite pair and bound the sum of products by the
product of the two total oriented masses, whose sum is $\mu_\ell$.
Summing over lines gives $P_-\le P_\parallel/2$ and $P_-\le1/2$.

Let $Y,Z>0$ bound the continuations with nonzero and zero sum,
respectively, uniformly over forms, previous zero equations, and
nonzero requirements. For $\tau\ge1$, an independent pair contributes
at most $Y+Z/\tau$ to the rank potential $\tau^{-\rank W}$: its zero
branch gains one rank. A dependent opposite pair contributes only $Z$;
a dependent nonopposite pair contributes only $Y$. In the latter case,
imposing zero sum would force both required nonzero inputs to vanish.

The with-replacement average is bounded by
\nopagebreak
\[
\begin{aligned}
 &(1-P_\parallel)(Y+Z/\tau)+P_- Z+(P_\parallel-P_-)Y\\
 &\quad=Y+Z/\tau+P_-(Z-Y)-P_\parallel Z/\tau\\
 &\quad\le Y+Z/\tau+P_-\bigl((1-2/\tau)Z-Y\bigr)\\
 &\quad\le\max\{Y+Z/\tau,(Y+Z)/2\}.
\end{aligned}
\]
The last step maximizes an affine function on $0\le P_-\le1/2$.
The labels are distinct with probability $(m-1)/m$. Since all branch
contributions are nonnegative, conditioning on this event gives
\begin{equation}\label{eq:fresh-pair-lemma}
 \begin{aligned}
 \mathbb E_{\rm distinct}[\text{branch contribution}]
 &\le\frac{m}{m-1}\Psi(Y,Z),\\
 \Psi(Y,Z)&=\max\left\{Y+\frac Z\tau,\frac{Y+Z}{2}\right\}.
 \end{aligned}
\end{equation}
The inequality also holds for $Y=0$ or $Z=0$ by continuity. In particular,
a forced zero sum costs at most
$\frac{m}{m-1}\max\{1/\tau,1/2\}$. The improvement comes from counting
proportional nonopposite pairs as dead zero branches, rather than allowing
them a rank increment. Further nonzero requirements can only remove branches.

For positive input and output markers $x,y$, define
\begin{align*}
 T_0(v)&=v_0+xyv_2,\\
 T_1(v)&=yv_1+x\max\{yv_1+v_0/\tau,v_0\},\\
 T_2(v)&=yv_2+x\Psi(yv_1,v_0).
\end{align*}
State 0 has zero feedback. State 1 has nonzero feedback after at least
two active inputs in the current run; $T_1$ is the nonzero-feedback
operator~\eqref{eq:width-one-quotient} at $q=1$.
State 2 is a \emph{macro state}: it stands for nonzero feedback after
exactly one active input in the current run, but $T_2$ bounds that input
together with the next active input through~\eqref{eq:fresh-pair-lemma}.
It does not describe a bound conditional on a revealed first label.
This distinction is essential: after revealing one label, every remaining
form might be its opposite.

To justify the macro state, write $A_t=(T^t\mathbf1)_0$ and
$B_t=(T^t\mathbf1)_1$. Expanding the state-2 recurrence gives
\[
 A_t=A_{t-1}+xy^t+
 \sum_{\gamma=0}^{t-2}x^2y^{\gamma+1}
 \Psi(yB_{t-\gamma-2},A_{t-\gamma-2}).
\]
The first term has zero input. The second is a final run with one
active input. In each summand, choose the gap $\gamma$ of zero inputs
before exposing either of the first two active labels.
Apply~\eqref{eq:fresh-pair-lemma}, then the uniform continuation bound.
State 1 uses the two-state bound of Section~\ref{sec:distance-proof}
until the feedback returns to zero. This proves the macro interpretation
without conditioning on the first label alone.

Each completed pair consumes two labels. The values $m$ at pair starts
are distinct and decrease along a trajectory, even if single exposures
intervene. Their sampling corrections therefore multiply to at most
$\prod_{m=2}^s m/(m-1)=s$ for a scan with input weight $s$.
More formally, the induction budget $f(m)=\max(1,m)$ satisfies both
$f(m-1)\le f(m)$ and $mf(m-2)/(m-1)\le f(m)$.
Final single-input runs require no pair correction.

Normalize $v_0=1$, choose $\lambda\ge\max_iT_i(v)/v_i$, and put
$C=\max_i1/v_i$. Monotonicity and homogeneity imply
$(T^n\mathbf1)_0\le C\lambda^n$. With the pair corrections, a scan with
$m$ active inputs therefore has the uniform bound $K(m)=mC\lambda^n$
required by Lemma~\ref{lem:scan-composition}. For source weight $s$ and
intermediate weight $h$, the lemma gives
\[
 \frac{F_s(\tau)}{\tau}
 \frac{sh\,C_1\lambda_1^nC_2\lambda_2^n}
 {\binom ns\binom nh\,x_1^s(y_1x_2)^h y_2^e},
 \qquad 0<y_2\le1.
\]
The lemma keeps the same zero-equation span and rank potential across
both scans. Conditioning on the first trajectory leaves the second
permutation independent and uniform, and the uniformity of the
continuation bounds justifies applying the macro state to its
intermediate forms.

For $e=209715$, the sparse bound~\eqref{eq:all-one-sparse} of
Appendix~\ref{app:all-one} bounds source weights 1 through 31 by a
probability whose base-two logarithm is at most $-47.919090471278$.
The dense range $32\le s\le4e$, $\lceil s/2\rceil\le h\le2e$ is
partitioned into 121 leaves, as in Appendix~\ref{app:all-one}. Each
leaf bounds $sh$ by the product of its rectangle's upper endpoints and
uses the rational binomial and convex vertex bounds of that appendix.
Exact rational evaluation, including all finite prefactors and
upward-rounded leaf sums, gives total logarithmic upper bound
$-47.919090461115$. The certificate and its verifier are listed in
Appendix~\ref{app:reproduction}. This strengthens the proof without
changing the encoder's $6.8125n-2$ field additions.
